%%%% ABSTRACT

\begin{abstractutfpr}%% Ambiente abstractutfpr
Public transport is a vital component for mobility in smart cities, generating large volumes of geolocation data that can be leveraged to improve operational management. However, the direct use of raw GPS data faces significant challenges due to noise, inaccuracies, and communication failures, hindering the faithful reconstruction of itineraries. This work proposes a methodology to apply Process Mining techniques to data from the Integrated Transport Network (RIT) of Curitiba (URBS). The main objective is to transform unstructured GPS records into a reliable event log, enabling the discovery and analysis of real operational processes. The methodology adopts the PM4Py library for mining and integrates Map Matching algorithms for the precise identification of bus stops, as well as temporal interpolation techniques to correct data gaps. As a result, a structured XES file was generated, enabling the application of discovery algorithms (Inductive Miner) and conformance checking. The analysis allowed the identification of performance bottlenecks and deviations between executed and planned routes, demonstrating the viability of spatial data engineering as a prerequisite for transport intelligence.
\end{abstractutfpr}
